\documentclass[11pt,a4paper]{article}
\usepackage[T1]{fontenc}
\usepackage[utf8]{inputenc}
\usepackage{lmodern}
\usepackage[margin=2.5cm]{geometry}
\usepackage{microtype}
\usepackage{amsmath}
\usepackage{enumitem}
\usepackage[hidelinks]{hyperref}
\setlist{itemsep=1pt,topsep=3pt,parsep=0pt}

\begin{document}
\begin{center}
  {\Large\bfseries Sonifying the Heart: Inner Emotion Detection from ECG through Music}\\[6pt]
  Music Workshop Sprint 3: Individual Contribution\\[8pt]
  {\large Sathwik Reddy} \quad \texttt{2024121002}
\end{center}
\noindent\rule{\textwidth}{0.4pt}

\section{My Role}
I built the fourth stage of our pipeline, \textbf{Emotion}, which reads the emotion of the
generated music. I also evaluated it.

\section{What I Did}
\begin{itemize}
  \item \textbf{Our emotion model.}
    \begin{itemize}
      \item It places each melody on Russell's valence--arousal plane, using three musical cues:
        rhythmic irregularity, pitch spread and consonance.
      \item Tempo and loudness are left out on purpose. Tempo differs more between people than
        between rhythms, and loudness is re-scaled within each clip.
    \end{itemize}
  \item \textbf{Essentia.}
    \begin{itemize}
      \item Four pretrained mood classifiers and a valence/arousal model trained on DEAM, run on
        both the melody and the full arrangement.
      \item Found and fixed a bug: the classifiers list their classes in different orders
        (\texttt{happy, non\_happy} but \texttt{non\_sad, sad}).
    \end{itemize}
  \item \textbf{music2emo.} Set up its separate environment and ran it on all 14 melodies.
  \item \textbf{Comparison.} For every recording, checked whether each method scores the AFib
    piece higher or lower than the normal one.
  \item \textbf{Validation plan.}
    \begin{itemize}
      \item ECG datasets with self-reported emotion (WESAD, DREAMER, AMIGOS), compared against
        plain heart-rate variability.
      \item The ethics of using this with specific groups, such as autistic people.
    \end{itemize}
\end{itemize}

\section{Key Result}
Our model reads all 7 normal pieces as calm, and 6 of the 7 AFib pieces as tense.

All three methods hear a difference in every recording, but no two agree on both valence and
arousal:
\begin{itemize}
  \item Essentia hears AFib as more pleasant and more energetic.
  \item music2emo hears it as less of both.
\end{itemize}
Tempo probably explains part of this, since every AFib clip is also the faster heart.

\section{Next Steps}
\begin{itemize}
  \item Run the validation on WESAD first.
  \item Design the listening study.
\end{itemize}

\end{document}
